\scanpage{51}
faut donc choisir un point \inserted{de $D$} au dessus
de $-\bar j$ (qui \unsure{permettra} alors d'identifier
$D$ au rev. universel \unsure{ponctué|pointé} en $P$ de
$U_{0,3}$) -- ce qui introduit justement
l'indétermination par un élément du
noyau de $\mathrm{SL}(\unsure{2|1},\mathbb Z)'\longrightarrow
\mathrm{SL}(2,\mathbb Z/2\mathbb Z)\simeq\unsure{\mathfrak S|\mathcal G}_3$,
lui-$\overline{\mathrm m}$ isomorphe à $\Pi_{0,3}$. Mais
on peut prendre p. ex.
\[
\tag{77}
j=\exp\frac{2i\pi}{3}=\frac{-1+i\sqrt{3}}{2}
\]
comme \deleted{\unsure{représentant}} pt. en dessus de
$P=-\bar j$ -- en tordant des \unsure{bras} pour
qu'il soit bien en dessus de $P$, et non
de $\bar P=-j$, sinon c'est $-\bar j=\dfrac{1+i\sqrt{3}}{2}=\exp\dfrac{2i\pi}{6}$
qu'on \unsure{désirerait|choisirait}. [\uline{NB} On trouve que
les trois éléments de $\mathrm{SL}(2,\mathbb Z)'\simeq\mathrm{GP}(1,\mathbb Z)^+$
qui transforment $j$ en $-\bar j$ sont
$\begin{pmatrix}1&1\\0&1\end{pmatrix}$,
$\begin{pmatrix}1&0\\\unsure{1}&1\end{pmatrix}$,
$\begin{pmatrix}0&-1\\\unsure{1|-1}&0\end{pmatrix}$, dont aucun
n'est $\equiv\begin{pmatrix}1&0\\0&1\end{pmatrix}\bmod 2$, ce qui montre que
$j$ et $-\bar j$ \deleted{ont des images \illegible{1}} ne sont pas congrus

\scanpage{52}
\mspage{314}
mod le noyau \inserted{$\Pi(2)$} de
$\mathrm{SL}(\unsure{2|1},\mathbb Z)\longrightarrow\mathrm{SL}(2,\mathbb Z/2\mathbb Z)$,
i.e. qu'ils n'ont pas $\overline{\mathrm m}$ image dans $U_{0,3}^{\mathrm{an}}
\simeq D/\Pi(2)$, donc leurs images sont
$P=-\bar j$ et $\bar P=-j$ (ou inversement)] Le
choix en question est \deleted{alors} \unsure{caractérisé}
par le fait que $\tau\in D$ en dessus de
$P$ est \deleted{\unsure{donc}} une \uline{racine de l'unité}, ce qui
est quand $\overline{\mathrm m}$ bien naturel. On peut
donc dire qu'on a ``normalisé'' l'hom
\inserted{ext.} (73) en un hom effectif -- et du
$\overline{\mathrm m}$ coup aussi (70).

\figure{[Court trait oblique descendant de gauche à droite, entre les paragraphes.]}

On aurait envie de reprendre les
constructions précédentes dans le contexte
des groupes profinis -- malheureusement
je manque des fondements nécessaires,
et surtout d'une théorie du ``\unsure{trouage}''
ou ``\unsure{forage} de \unsure{trous}'' au niveau groupes
à lacets profinis -- même de \deleted{\illegible{2}} \inserted{\unsure{façon}} \unsure{conjec-}

\scanpage{53}
\unsure{turale}. Pourtant on sent que c'est
bien dans cette direction que se trouve
une compréhension ``\unsure{gruppentheoretisch}''
de l'homomorphisme \deleted{$\unsure{\mathfrak S|\mathcal G}_{0,3}$} (40)
\[
\widehat{\unsure{\mathcal M|M}}_{0,3}\longrightarrow
\mathrm{GL}(2,\widehat{\mathbb Z})/\mathord{\pm}1
\simeq\mathrm{GP}(1,\widehat{\mathbb Z}),
\]
qui est \unsure{la|le} motivation principale des
réflexions du présent \S, et des \unsure{précédents}.
Ce devrait être une incitation pour une
réflexion \unsure{renouvelée} sur la question du
\unsure{trouage} dans le contexte profini\ldots Mais
je me rends compte dès maintenant
que \deleted{\unsure{l'}} \inserted{la question s'il y a une} extension des constructions du
présent \S, au contexte profini, indépendamment
de l'introduction d'éléments
de structure ``arithmétique'' (tel que
\unsure{la|le} \uline{\unsure{donnée|donné}} d'un sous-groupe
$\Gamma\subset\mathrm{Autext}_{\mathrm{loc}}(\widehat{\Pi}_{0,3}\unsure{)}$,
satisfaisant à un certain nombre de
conditions $\pm$ subtiles, genre relations
congruentielles) est \unsure{subordonnée} à celle,

\scanpage{54}
\mspage{315}

si ces relations congruentielles sont bel et
bien des relations de nature arithmétique,
spécifiques aux autom. ext. \unsure{provenant}
de $\Gamma_Q$, ou au contraire des
conséquences $\pm$ automatiques d'un
formalisme \inserted{\unsure{purement} ``algébrique''} des groupes profinis, \inserted{qui serait} \unsure{calqué}
sur le cas discret. Sans préjuger \unsure{aucunement}
de la réponse et du degré de
plausibilité d'un \unsure{tel} formalisme, je
pourrais à titre expérimental essayer
d'en dégager un, pour \unsure{tester|tenter} où il
me mènera.

\uline{\S 41} Condition de commutation de
$\widehat{\Pi}'_{0,3}\xrightarrow{\sim}\widehat{\unsure{\mathfrak S|\mathcal G}}_{0,3}^{+}$
aux opérations de $\Gamma_Q$ : Calcul d'autom.
extérieurs de $\widehat{\Pi}_{0,3}$ et $\widehat{\unsure{\mathfrak S|\mathcal G}}_{0,3}^{+}$.

Avant de me lancer dans cette
\deleted{\unsure{tentative}}\inserted{voie}, je voudrais encore élucider
une question, concernant d'autres
éventuelles conditions nécessaires (en plus
des conditions congruentielles du \S\unsure{38|39})
pour qu'un autom. ext. de $\widehat{\Pi}_{0,3}$ provienne

\scanpage{55}

de $\Gamma_Q$. Pour \unsure{ce|ceci}, reprenons l'hom
extérieur composé
\[
(1)\qquad
\diagram{\begin{tikzcd}[ampersand replacement=\&,column sep=small,row sep=small]
\Pi_{0,3} \arrow[r,"\text{épi}"] \arrow[rr,bend left=25,"\varphi"]
\& \Pi'_{0,3} \arrow[r,"\sim"]
\& \unsure{\mathfrak S|\mathcal G}_{0,3}^{+}
  \simeq \unsure{\pi_1|\Pi_1}\bigl((U_{0,3}^{\mathrm{an}},\unsure{\mathfrak S|\mathcal G}_3),\unsure{P}
\\
\& {}=\Pi_{0,3}/\{\ell_1^2,\ell_\infty^3\}
\\
\& {}\simeq\text{groupe cartographique}
\\
\& \text{\unsure{triangulé} orienté}
\end{tikzcd}}
\]
défini par
\[
(2)\qquad
\left\{\begin{aligned}
\varphi(\ell_0)&=\ell'_0=\sigma_\infty\rho\\
\varphi(\ell_1)&=\ell'_1=\sigma_\infty\\
\varphi(\ell_\infty)&=\ell'_\infty=\rho^{-1}
\end{aligned}\right.
\]
(avec les notations de [C, \ldots] \unsure{\ldots}),
qui par passage aux complétés donne
un hom extérieur
\[
(3)\qquad
\diagram{\begin{tikzcd}[ampersand replacement=\&,column sep=small,row sep=small]
\widehat{\varphi}:\widehat{\Pi}_{0,3} \arrow[r,"\text{épi}"]
\& \widehat{\unsure{\mathfrak S|\mathcal G}}_{0,3}^{+} \arrow[r,phantom,"\simeq"]
\& \unsure{\pi_1|\Pi_1}\bigl((\unsure{U_{0,3\,\overline Q_0}},\unsure{\mathfrak S|\mathcal G}_3),P\bigr).
\\
\& \cap
\\
\& \unsure{\widehat{\widehat{\mathcal T}}}_{0,3} \arrow[r,phantom,"="]
\& \operatorname{Aut}\deleted{\unsure{\mathrm{int}|\mathrm{ext}}}_{\mathrm{loc}}(\widehat{\Pi}_{0,3})
\end{tikzcd}}
\]

Cet homomorphisme est défini ici \deleted{\illegible{1}}\inserted{par les}
``formules brutales'' (2), sans
référence à un sous-groupe particulier
$\Gamma$ de $\operatorname{Aut\,ext}'_{\mathrm{loc}}(\widehat{\Pi}_{0,3})$. Si on se donne
un \deleted{\illegible{1}} tel sous-groupe, commutant
à $\unsure{\mathfrak S|\mathcal G}_3$, celui-ci opère automatiquement
sur le groupe extérieur $\widehat{\unsure{\mathfrak S|\mathcal G}}_{0,3}^{+}$, car il

\scanpage{56}
\mspage{316}

s'insère dans une suite \unsure{exacte|ex.te}
\[
(4)\qquad
\diagram{\begin{tikzcd}[ampersand replacement=\&]
1\arrow[r]
\&\widehat{\unsure{\mathfrak S|\mathcal G}}_{0,3}^{+}\arrow[r]
\&M_{0,3}\arrow[r]
\&\Gamma\arrow[r,dash]
\&1.
\end{tikzcd}}
\]
\deleted{en plus}\inserted{\unsure{contenant}} la suite \unsure{exacte|ex.te}

(5)\quad
\diagram{[Deux rangées. En haut, de gauche à droite :
$1$, $\widehat{\Pi}_{0,3}$, $\mathcal E_{0,3}$, $\Gamma$, $1$.
En bas, de gauche à droite :
$1$, $\widehat{\unsure{\mathfrak S|\mathcal G}}_{0,3}^{+}$, $M_{0,3}$, $\Gamma$, $1$.
Chaque paire d'objets voisins d'une même rangée est reliée par un trait horizontal ;
les pointes discernables vont vers la droite, mais celles des derniers traits sont peu distinctes.
De $\widehat{\Pi}_{0,3}$ vers le groupe situé au-dessous : trait vertical à sommet recourbé,
terminé par une pointe vers le bas.
Entre $\mathcal E_{0,3}$ et $M_{0,3}$ : trait vertical recourbé au sommet, sans pointe inférieure discernable.
Entre les deux $\Gamma$ : deux petits traits parallèles verticaux.
Un point suit la rangée inférieure.]}

Le fait que l'homomorphisme
(3) commute à l'action de $\Gamma_Q=\Gamma$
me semble \deleted{\unsure{être}} a priori une condition
non triviale sur $\Gamma$ -- plus
précisément, sur les éléments de
$\Gamma$ séparément, qu'il faudrait bien
\unsure{arriver} à expliciter (de $\overline{\mathrm m}$ que les
conditions congruentielles), en
termes des calculs $\pm$ explicites
de [C]. Cette condition équivaut
à celle que l'hom (3) s'insère
dans un homomorphisme d'extensions

\scanpage{57}

\margin{NB. $\mathcal E_{0,3}$ est le groupe noté $M_{0,3}$! dans le \S\,\unsure{46}}

\diagram{\[
(6)\qquad
\begin{tikzcd}[ampersand replacement=\&]
1 \arrow[r] \& \widehat{\Pi}_{0,3} \arrow[r] \arrow[d,"\widehat{\varphi}"] \& \mathcal E_{0,3} \arrow[r] \arrow[d,"\Phi"] \& \Gamma \arrow[r] \arrow[d,"1"] \& 1 \\
1 \arrow[r] \& \widehat{\unsure{\mathfrak S|\mathcal G}}_{0,3}^{+} \arrow[r,dash] \& M_{0,3} \arrow[r] \& \Gamma \arrow[r,dash] \& 1
\end{tikzcd}
\]}
,

\deleted{qui s'identifie à}
moyennant lequel $M_{0,3}$ (\unsure{interprété} comme groupe fondamental ``à signature'' de $(\mathbb P^1_{\unsure{\mathbb Q|Q}}-\{0\},\,2\{1\}+3\{\infty\})$) s'identifie au quotient de $\widehat{\Pi}_{0,3}$ par le sous-groupe invariant \inserted{fermé} engendré par $\ell_1^2,\ell_\infty^3$.

Je commence à \unsure{soupçonner} qu'il y aura toute une \unsure{kyrielle} de conditions de même nature sur $\Gamma$, qu'il faudrait bien arriver à dégager de façon plus générale.

\scanpage{58}
\mspage{317}

\unsure{Je|je} vais reprendre les calculs faits dans [C], concernant les \unsure{autom.|autom} extérieurs de $\widehat{\Pi}_{0,3}$ qui commutent à $\unsure{\mathfrak S|\mathcal G}_3$, \deleted{\unsure{qui|que}} \inserted{Ils \deleted{peuvent}} peuvent se \unsure{réaliser} de façon unique comme \unsure{automorphismes} \unsure{effectifs} \inserted{$u$} qui \unsure{commutant|commutent} à \unsure{$\sigma_\infty$;|$\sigma_\infty'$,} (et qui de plus \unsure{commutent|commutant} extérieurement à $\rho$). Ils correspondent \inserted{biunivoquement} aux couples

\unsure{A.}
\[
(7)\qquad (g_0,p)\qquad
\left\{
\begin{aligned}
g_0&\in\widehat{\Pi}_{0,3}\bmod L_0\\
p&\in\widehat{\mathbb Z}^{*}\quad (\text{multiplicateur})
\end{aligned}
\right.
\]
satisfaisant les conditions suivantes (où $\deleted{\unsure{\beta_0|\rho_0}}\,\unsure{\sigma|\gamma}\in\widehat{\mathbb Z}$ \deleted{\unsure{doit être}} \inserted{est} déterminé de façon unique par $u$ (ou encore, par l'autom. extérieur correspondant)), où on a posé
\[
h_0=\sigma_\infty(g_0)^{-1}g_0=[\sigma_\infty,g_0^{-1}]:
\]

\margin{\uline{NB} \unsure{Admettant} que $\widehat{\Pi}_{0,3}^{\sigma_\infty}=\{1\}$, on voit que $g_0$ est \unsure{(?)} \unsure{connu} \unsure{quand} on connaît $h_0$\ldots}

\margin{\uline{NB} en fait, \unsure{on} \unsure{verra|voit} que (9) \unsure{est} \unsure{superflu|superflue}, et $\unsure{\sigma|\gamma}=\frac{p-1}{2}$.}

\[
\begin{array}{c@{\quad}l}
\begin{array}{c}(9)\\[1em](10)\end{array}
&
\boxed{\begin{aligned}
&\ell_0^p\cdot\operatorname{int}(h_0)(\ell_\infty^p)\operatorname{int}(\ell_1^{\unsure{\sigma|\gamma}}\rho(h_0)^{-1})(\ell_\infty^p)=1\\[.5em]
&h_0\ell_0^{-\unsure{\sigma|\gamma}}\rho^{-1}(h_0)=\ell_1^{\unsure{\sigma|\gamma}}\rho(h_0)^{-1}\ell_\infty^{\unsure{\sigma|\gamma}}
\end{aligned}}
\end{array}
\]

\deleted{\illegible{2} que} L'hom $u$ est donné par

\figure{[Deux longs traits presque parallèles sous la dernière ligne de texte, au-dessus de la note marginale inférieure.]}

\margin{\uline{NB} (10) s'écrit aussi \deleted{\unsure{(11)}}
\[
(10\text{ bis})\qquad
(\ell_{\unsure{\infty|0}}^{\unsure{\sigma|\gamma}}h_0^{-1})\cdot\rho(\ell_{\unsure{\infty|0}}^{\unsure{\sigma|\gamma}}h_0^{-1})\cdot\rho^2(\ell_{\unsure{\infty|0}}^{\unsure{\sigma|\gamma}}h_0^{-1})=1
\]
}

\scanpage{59}
\mspage{318}

\[
(11)\qquad
\left\{
\begin{aligned}
u(\ell_0)&=\operatorname{int}(g_0)(\ell_0^p)\\
u(\ell_1)&=\operatorname{int}\bigl(\underbrace{\sigma_\infty(g_0)}_{\inserted{\unsure{g_1}}}\bigr)(\ell_1^p)\\
u(\ell_\infty)&=\operatorname{int}\bigl(g_0\ell_\infty^{-\unsure{\sigma|\tau}}\rho^{-1}(h_0)\bigr)(\ell_\infty^p)
\end{aligned}
\right.,
\]

et il faut en toute rigueur, en plus des conditions (9),(10), exiger que les deux éléments $u(\ell_0),u(\ell_1)$ \inserted{donnés par (11)} \unsure{forment} des générateurs du groupe profini $\widehat{\Pi}_{0,3}$.
\margin{(11 bis)}

Je vais calculer l'action de $u$ sur $\widehat{\unsure{\mathfrak S|\mathcal G}}_{0,3}^{+}$, en \unsure{considérant} ce groupe comme engendré par les éléments
\[
(12)\qquad \ell'_1=\sigma_\infty,\qquad \ell'_\infty=\rho^{-1}
\]
donc il faut calculer simplement \inserted{$\operatorname{int}(u).\sigma_\infty=u(\sigma_\infty)=$} $u\sigma_\infty u^{-1}$ et $u(\rho)=\operatorname{int}(u).\rho=u\rho u^{-1}$. On a par construction ($u$ \unsure{commutant} à $\sigma_\infty$)
\[
(13)\qquad u\unsure{\circ}\sigma_\infty u^{-1}=\sigma_\infty,
\]
et d'autre part on vérifie que
\[
\deleted{u(\rho)=u\rho u^{-1}=\operatorname{int}(h}
\]
\[
(14)\qquad \deleted{u\rho}\quad
\left\{
\begin{aligned}
\rho^{-1}u\rho&=\operatorname{int}(h)u\\
\inserted{\text{avec}}\quad h&=(\rho^{-1}\sigma_\infty)(g_0)\ell_\infty^{\unsure{\gamma|\sigma}}g_0^{-1}
\end{aligned}
\right.
\]
\margin{i.e. $u\rho u^{-1}=\rho\operatorname{int}(h).$\\
$=\operatorname{int}(\rho(h))\rho$\\
i.e. $\rho(h)=\sigma_\infty(g_0)\ell_1^{\unsure{\gamma|\sigma}}\rho(g_0)^{-1}$}

donc on trouve, en \unsure{termes} d'une opération sur le groupe $\widehat{\unsure{\mathfrak S|\mathcal G}}_{0,3}^{+}$ engendré par \inserted{$\ell'_1=$}$\sigma_\infty$ et \inserted{$\ell'_\infty=$}$\rho^{-1}$ (ou $\rho$)
\[
(15)\qquad
\left\{
\begin{aligned}
u(\sigma_\infty)&=\sigma_\infty\\
u(\rho)&=g.\rho
\end{aligned}
\right.
\quad\text{avec}\quad
 g=\sigma_\infty(g_0)\ell_1^{\unsure{\tau|\sigma}}\rho(g_0)^{-1}
\]

\scanpage{60}

\uline{NB} La relation \deleted{$u(\rho)^3=1$} $u(\sigma_\infty)^2=1$ est triviale, la relation $u(\rho)^3=1$ se réduit à
\[
\deleted{g\rho(g)\unsure{=}}
\]
\[
(16)\qquad g\rho(g)\rho^2(g)=1
\]
qui est conséquence de l'expression explicite (15) de $g$, et de la formule (10).

Notons qu'on \deleted{\unsure{aura}} \inserted{devrait avoir}, \unsure{\&|en} \unsure{reprenant}
\[
\ell'_0=\sigma_\infty\rho=\unsure{\varepsilon_0|e_0}\qquad
(\text{\unsure{d'où}}\quad \ell_0^{\prime 2}=\ell_0)
\]
\[
(17)\qquad
u(\unsure{\ell'_0})
=\operatorname{int}(g_0)(\ell_0^{\prime p})
\]
\inserted{\figure{[Sous l'argument de $u$, une petite flèche descendante mène à un signe entouré, $\unsure{\varepsilon}$.]}}
\margin{?}
--- cette relation \unsure{au|en} carré n'est autre que la première relation (11)
\[
u(\ell_0)=\operatorname{int}(g_0)(\ell_0^p).
\]

[La relation (17) en est une conséquence, si on admet que le centralisateur de $L_0=\ell_0^{\widehat{\mathbb Z}}$ dans $\widehat{\Pi}_{0,3}$ est réduit à $L_0$\dots]

On devrait avoir aussi la relation
\[
(18)\qquad u(\rho)\deleted{\unsure{\equiv|=}}\inserted{\text{conjugué de}}\rho^p
\quad\deleted{\text{\unsure{mod|modulo}}\ \widehat{\Pi}_{0,3}}
\]
\deleted{\illegible{1}} sur laquelle nous \unsure{allons} revenir -- il n'est pas clair pour moi qu'elle soit conséquence des conditions \unsure{envisagées} sur $u$ jusqu'à présent.
\margin{\deleted{\illegible{2}}}
